\section{Conclusion} \label{sec:conclusion}

This paper evaluated proficiency gains associated with adoption of the Full-Time Education Program in São Paulo state schools using SARESP data from 2011 to 2018. The central question was not only whether PEI schools improved after conversion, but whether this gain remains once the empirical design is taken seriously: staggered adoption, heterogeneous effects over time, observable differences across schools, and compositional change in the student body. The main estimates indicate average post-treatment effects between approximately 0.50 and 0.68 standard deviations, with larger effects in Mathematics and in high school. The increasing dynamics in the event studies could reflect dynamic student selection, changes in exam participation, or recovery among schools initially below the mean; they are also compatible with pedagogical and organizational maturation of the program. The available data do not allow these channels to be fully separated.

The substantive contribution of the paper is to show that these results do not depend only on a simple comparison between treated schools and controls. The gains remain when TWFE is abandoned as the main estimator, when \textcite{callawaysantanna} is used to handle staggered adoption, when covariates are added through \textcite{santannazhao}, and when an aggregate approximation inspired by \textcite{santannaxu} is implemented. This last step is important because the literature on PEI shows that school conversion changes student composition. Therefore, an evaluation of the program that ignores this channel risks confusing learning gains with student selection. Within the observable dimensions available in this paper, the compositional correction does not eliminate the estimated effects, but the Hausman-type diagnostic has low power and is not evidence against compositional effects of up to approximately 0.20--0.25 standard deviations.

The methodological contribution is to make explicit the distance between the ideal estimator and the feasible estimator. The estimator of \textcite{santannaxu} was designed for repeated cross-sections in which compositional change can be handled at the individual level. To apply it fully, one would need to observe, in the same record, individual enrollment, student covariates, and SARESP scores. This structure is not publicly available for the period analyzed. For this reason, this paper builds an aggregate adaptation: for each cohort and relative time, I estimate 2x2 cells using a four-category multinomial logit, composition weights, and outcome regressions, always using covariates fixed at the cohort baseline. The final version does not use a binary S\&Z fallback; when common support is insufficient, the cell is excluded. This makes the estimate more restrictive, but also more honest.

This limitation defines the scope of interpretation. The results are evidence of aggregate proficiency gains in PEI schools; they are not a perfect decomposition between individual learning, student selection, and changes in exam participation. Even so, the exercise is informative. If observable selection were the main explanation for the gains, one would expect the estimates to fall substantially when moving from CS21 to DR-DiD and to the S\&X approximation. This is not what occurs. In addition, the diagnostics for mean reversion and regional spillovers preserve magnitudes close to the main estimates. On the other hand, the no-imputation sensitivity check shows that the magnitude depends on coverage of enrollment covariates, especially because the complete-case sample is much smaller and, in high school, insufficient to sustain the same estimation. The participation diagnostic for SARESP also recommends additional caution in high school, where the ratio between participants and enrollment falls after conversion.

The main open questions follow from this same data restriction. The first natural extension is to obtain access to administrative databases that make it possible to link individual enrollment, exam participation, and individual scores, approximating the Sant'Anna and Xu estimator more faithfully and separating learning, selection, and exam attendance more clearly. The second is to incorporate more direct geographic measures of exposure to nearby PEI schools, using distance or commuting time, to move beyond the proxy based on school district and neighborhood. The third is to extend the window beyond 2018, separately treating the pandemic, changes in assessment, and the accelerated expansion of PEI. In sum, this paper shows that PEI is associated with persistent proficiency gains in an aggregate school panel. The conclusion does not settle the discussion about mechanisms, but it organizes its terms more clearly: there is a positive aggregate gain, there is compositional change, there is evidence of altered high-school participation, and the observable part of composition does not appear sufficient to overturn the results under the estimated specifications. At the same time, the magnitudes obtained here are larger than those estimated by \textcite{scorzafave2023} for the same program, reinforcing that the prudent interpretation is to treat the results as aggregate evidence conditional on the available design, not as a definitive measure of the individual net effect of PEI.
